% Imporditud: tools/import_eksperimendid.py
% Allikas: Eesti füüsikaolümpiaadi eksperimendi ülesannete
%   kogu (Rael Kalda, 2023), ülesanne Ü83, lahendus L83.
\setAuthor{EFO žürii}
\setRound{lõppvoor}
\setYear{2000}
\setNumber{P E2}
\setDifficulty{3}
\setTopic{elektriahelad}
\setVahendid{Taskulambipirn alusel, lapik patarei, tester, juhtmed}
\setPunktid{12}
\setAllikas{Eksperimendi kogu Ü83, lahendus lk 69}
\setLahendusseis{toores}

\prob{Taskulambipirni takistus}
Leida põleva taskulambipirni takistus.
\textit{Märkus:} i: Kasutatava testri takistus ampermeetrina mõõtepiirkonnal 10 A on 0,2 $\Omega$ ja mõõtepiirkonnal 200 mA --- 6 $\Omega$. Testri kasutamisel voltmeetrina mõõtepiirkondadel 0,02 V ja 2000 mVon takistus 1 M$\Omega$. NB! Testri kui oommeetri kasutamine pin- gestatud detailide takistuse määramiseks on keelatud, kuna see võib põhjustada testri rik- nemise.

\hint

\solu
Ühendame jadamisi patarei, taskulambipirni ja testri ampermeetrina mõõtepiirkonnal 10 A. Registreerime voolutugevuse I (2 p). Seejärel ühendame lambi otse patareiga ning testri voltmeetrina rööbiti lambiga. Registreerime klemmipinge U (2 p). Arvutame lambi takistuse eeldusel, et ampermeeter ja voltmeeter on ideaalsed: R = U/I (2 p). Võrreldes tulemust mõõteriistade takistustega, veendume selles, et tehtud eeldus on piisavalt hästi põhjendatud (2 p). Täheldame, et voolu mõõtmine piirkonnal 10 A on väga ebatäpne (1 p). Seetõttu teostame samad mõõtmised ka ampermeetri mõõtepiirkonnal 200 mA (1 p). Eeldades klemmipinge konstantsust, arvutame lambi ja ampermeetri jadaühenduse kogutakistuse Rk = U/I (1 p). Lambi takistuse R leidmiseks lahutame kogutakistusest ampermeetri takistuse. Seega R = Rk $-$ RA. (1 p). Kokku 12 p.

12 p võib ka teenida, rakendades keskkooli programmi kuuluvaid teadmisi: 1) mõõdame elektromotoorjõu E kui koormamata patarei klemmipinge (3 p); 2) mõõdame piirkonnal 10 A patarei lühisvoolu Il (3 p); 3) arvutame patarei sisetakistuse

r = (E/Il) $-$ RA (2 p);

4) mõõdame voolutugevuse I läbi lambi ja ampermeetri jadaühenduse (2 p); 5) arvutame lambi takistuse valemist

R = (E/I) $-$ r $-$ RA (2 p).
\probend
